\section{Continuous Microlocal Extension}
\label{sec:continuous}

Why should a learned packet layer follow a transported-symbol law at all? The answer comes from classical microlocal analysis. Short-time hyperbolic propagators are Fourier integral operators associated with canonical graphs, and their packet matrices are concentrated near those graphs. This section turns that classical fact into the continuous bridge used by the paper.

The finite landing spine is machine-checked. The continuous layer explains why sparse canonical packet matrices are natural for a concrete PDE family. The short-time wave propagator furnishes the first main continuous theorem, which in turn feeds an abstract canonical-graph packet sparsity theorem. The result is a dual spine:
\begin{enumerate}
\item a \emph{short-time half-wave off-graph decay theorem};
\item an \emph{abstract canonical-graph packet sparsity theorem}.
\end{enumerate}
Together they explain why the theorem-grade \MiNO\ class should use a bounded canonical stencil rather than either a dense interaction map or a purely single-destination law.

\begin{figure}[t]
\centering
\resizebox{\textwidth}{!}{%
\begin{tikzpicture}[>=Latex,font=\small,node distance=0.75cm]
\tikzstyle{box}=[draw,rounded corners,thick,align=center,minimum height=0.78cm]
\tikzstyle{classical}=[box,fill=blue!6,minimum width=3.25cm]
\tikzstyle{finite}=[box,fill=green!7,minimum width=3.25cm]
\tikzstyle{learned}=[box,fill=orange!9,minimum width=3.25cm]
\node[classical] (fio) {classical FIO input\\half-wave parametrix};
\node[classical,right=0.55cm of fio] (offgraph) {packet off-graph\\decay};
\node[classical,right=0.55cm of offgraph] (sparsity) {canonical-tube\\packet sparsity};
\node[finite,below=0.85cm of offgraph] (finite) {Lean-checked finite landing\\shells, tails, stencil sums};
\node[learned,right=0.55cm of sparsity] (net) {network-realizable\\transport and symbol budgets};
\node[box,fill=black!5,below=0.85cm of sparsity,minimum width=4.15cm] (mpt) {Microlocal Propagation Theorem\\for \MiNO-Core};
\draw[->,thick] (fio) -- (offgraph);
\draw[->,thick] (offgraph) -- (sparsity);
\draw[->,thick] (sparsity) -- (net);
\draw[->,thick] (sparsity) -- (finite);
\draw[->,thick] (finite) -- (mpt);
\draw[->,thick] (net) -- (mpt);
\node[align=center,above=0.08cm of offgraph,text=blue!55!black] {classical continuous bridge};
\node[align=center,left=0.08cm of finite,text=green!40!black] {machine-checked\\finite bridge};
\node[align=center,above=0.08cm of net,text=orange!70!black] {standard NN\\approximation input};
\end{tikzpicture}
}
\caption{Proof-boundary map for the continuous extension.  The oscillatory-integral and packet almost-diagonalization inputs are classical; the shell/tail/stencil landing layer is the machine-checked part; the learned transport and symbol terms enter through compact-window neural approximation budgets.}
\label{fig:proof-boundary-map}
\end{figure}

\subsection{Classical ingredients}

Write $S^m=S^m_{1,0}$ for the standard H\"ormander symbol class on $\R^d_x\times\R^d_\xi$, and $\Op(a)$ for the Kohn--Nirenberg quantization of $a\in S^m$. We use the following standard ingredients:
\begin{enumerate}
\item Calder\'on--Vaillancourt continuity for $\Op(a)$ with $a\in S^0$;
\item symbolic calculus for compositions of pseudodifferential operators;
\item short-time Fourier-integral parametrices for hyperbolic propagators generated by real principal symbols;
\item packet almost-diagonalization for canonical-graph Fourier integral operators;
\item non-stationary phase estimates obtained by repeated integration by parts.
\end{enumerate}
The relevant references are \citet{hormander1985analysis3,hormander1985analysis4,zworski2012semiclassical,alazardmicrolocal,candesdemanet2005curvelet}. The contribution of the present section is not to replace those results. It is to organize them around the exact packet objects used by \MiNO\ and to expose the proof chain in a form that can be mirrored by a Lean-friendly finite abstraction.

\subsection{Packet frames, shells, and canonical tubes}

Fix a semiclassical scale $h\in(0,h_0]$ and a wavepacket frame $\{\varphi_\gamma^h\}_{\gamma\in\Gamma_h}$ with packet centers
\[
z_\gamma=(x_\gamma,\xi_\gamma)\in T^\ast\R^d,
\]
analysis map $W_hu=(\langle u,\varphi_\gamma^h\rangle)_\gamma$, synthesis map $W_h^\ast$, and bounded overlap.  Let $d$ denote raw phase-space distance and set $d_h=h^{-1/2}d$.  Throughout this section $\delta_h$ is the \emph{normalized} covering defect, so the raw covering radius is $h^{1/2}\delta_h$.  For a canonical map $\kappa$ and a source packet $\eta$, we measure output packets relative to the propagated center $\kappa(z_\eta)$ using the shell index
\[
\ell_\kappa(\gamma,\eta)
:=
\Bigl\lfloor h^{-1/2}\,\mathrm{dist}\bigl(z_\gamma,\kappa(z_\eta)\bigr)\Bigr\rfloor.
\]
The corresponding canonical tube of radius parameter $\rho$ is the set of packets with shell index at most $\rho$. The finite Lean files introduced in this revision formalize exactly this shell-and-tube bookkeeping layer, independently of the continuous oscillatory-integral input.

\begin{assumption}[Quantitative packet-frame constants]
\label{ass:packet-frame-constants}
Throughout the continuous bridge, the packet family is restricted to a compact phase window $\Omega\subset T^\ast\R^d$ and satisfies the following scale-uniform bounds:
\begin{enumerate}
\item \emph{Frame stability:} there are constants $0<A_W\le B_W<\infty$, independent of $h$, such that
\[
A_W\|u\|_{L^2}^2
\le
\sum_{\gamma\in\Gamma_h}|\langle u,\varphi_\gamma^h\rangle|^2
\le
B_W\|u\|_{L^2}^2
\]
for all $u$ microlocalized in $\Omega$.
\item \emph{Covering:} the packet centers cover $\Omega$ at raw radius $h^{1/2}\delta_h$, with $\delta_h\le C_{\mathrm{cov}}$.
\item \emph{Counting:} canonical balls of radius $Rh^{1/2}$ contain at most $C_{\mathrm{cnt}}R^q$ packet centers.
\item \emph{Uniform localization:} packet envelopes and their derivatives up to the order used in stationary and non-stationary phase are bounded by constants collected in $C_{\mathrm{loc},N}$ after the standard $h^{1/2}$ rescaling.
\end{enumerate}
\end{assumption}

The constants in the continuous theorems below are allowed to depend on
\[
\mathfrak C_{\mathrm{frame},N}
:=
(A_W^{-1},B_W,C_{\mathrm{cov}},C_{\mathrm{cnt}},C_{\mathrm{loc},N})
\]
and on the compact-window symbol and flow seminorms. They are not allowed to depend on $h$, the number of packets, the retained stencil size $K$, or the branch parameter budgets. This convention is the quantitative version of the packetization defect tracked in the experiments: changing from local FFT packets to Gaussian/Gabor or multiscale Gabor packets changes $\mathfrak C_{\mathrm{frame},N}$ and $\delta_h$, hence it changes the theorem constants rather than the algebraic finite bridge.

\begin{lemma}[Constant ledger for the propagation theorem]
\label{lem:constant-ledger}
In Theorem~\ref{thm:network-realizable-wave-mino} and the Microlocal Propagation Theorem below, the constants may be chosen with the following dependency classes:
\[
C_{\mathrm{tail}}=C(C_{\mathrm{cnt}},q,N,C_{T,N}),
\qquad
C_{\mathrm{disc}}=C(T,M,\mathfrak C_{\mathrm{frame},M},\|a_t\|_{S^0,M}),
\]
\[
C_{\mathrm{tr}}=C(K_{\kappa,T},\mathfrak C_{\mathrm{frame},2},\|s_t\|_{C^1},C_{\mathrm{deg}}),
\qquad
C_{\mathrm{sym}}=C(B_W,C_{\mathrm{deg}},\|s_t\|_{L^\infty}),
\]
and
\[
C_{\mathrm{field}}=C(A_W^{-1},B_W,E^{\mathrm{in}}_{\mathrm{proj}},E^{\mathrm{out}}_{\mathrm{proj}}).
\]
Here $K_{\kappa,T}$ denotes compact-window Lipschitz and inverse-Lipschitz bounds for the no-caustic canonical graph, $\|a_t\|_{S^0,M}$ collects the amplitude seminorms used by stationary phase, and $C_{\mathrm{deg}}$ is the row-column stencil conversion constant from Lemma~\ref{lem:row-column-stencil}. These constants do not depend on $h$, the packet count, $K$, $P_{\mathrm{tr}}$, or $P_{\mathrm{sym}}$ once those quantities are kept in the explicit error terms.
\end{lemma}

\begin{proof}
The tail constant is the dyadic shell-sum constant from Theorem~\ref{thm:packet-sparsity}. The discretization constant is the stationary-phase and packet-covering constant from Theorem~\ref{thm:continuous-reduction}. The transport constant is the Lipschitz conversion from canonical-map error to packet-kernel mismatch on the retained tube, followed by the bounded-degree row-column transfer. The symbol constant is the retained-stencil multiplier bound. The field constant is the frame synthesis and projection constant used to pass from packet coefficients back to $L^2$ fields. Each step uses only compact-window seminorms and the frame constants listed in Assumption~\ref{ass:packet-frame-constants}.
\end{proof}

\subsection{Spine A: short-time half-wave off-graph decay}

We now specialize to the cleanest PDE family in the paper. Let $c\in C^\infty(\R^d)$ satisfy
\[
0<c_{\min}\le c(x)\le c_{\max}<\infty,
\]
and define the half-wave principal symbol
\[
p(x,\xi)=c(x)|\xi|.
\]
Let $U_+(t)$ denote the positive half-wave propagator generated microlocally by $p$, and let $\kappa_t$ be the Hamiltonian flow of $p$. We work on a time interval $|t|\le T$ chosen so that the flow remains a smooth canonical graph on the frequency band seen by the packet frame, with no caustic formation or branch crossing.

\begin{theorem}[Short-time half-wave off-graph decay]
\label{thm:wave-offgraph}
Fix $T>0$ inside a no-caustic window for the half-wave flow generated by $p(x,\xi)=c(x)|\xi|$. Let $\{\varphi_\gamma^h\}_{\gamma\in\Gamma_h}$ be an analytic Gaussian/Gabor-type packet frame with normalized covering defect $\delta_h$. Then for every integer $N\ge 1$ there is a constant $C_{T,N}$ such that
\begin{equation}
\label{eq:offgraph}
\bigl|\langle U_+(t)\varphi_\eta^h,\varphi_\gamma^h\rangle\bigr|
\le
C_{T,N}
\Bigl(1+h^{-1/2}\,\mathrm{dist}\bigl(z_\gamma,\kappa_t(z_\eta)\bigr)\Bigr)^{-N}
\end{equation}
uniformly for $|t|\le T$, $h\in(0,h_0]$, and $\gamma,\eta\in\Gamma_h$.
\end{theorem}

\begin{proof}[Proof sketch]
The proof has four classical steps.

\paragraph{Step 1: parametrix.}
For $|t|\le T$, the short-time half-wave propagator is microlocally represented by an oscillatory integral
\[
U_+(t)u(x)
\sim
(2\pi h)^{-d}\iint
e^{\frac{i}{h}(\Phi_t(x,\xi)-y\cdot\xi)}
a_t(x,\xi;h)\,u(y)\,dy\,d\xi,
\]
where $\Phi_t$ generates the canonical graph of $\kappa_t$ and $a_t$ is a classical amplitude. The specific input is the usual real-principal-type short-time parametrix: on a compact frequency band before caustic formation, the half-wave group is a semiclassical Fourier integral operator associated with the canonical graph of the Hamiltonian flow. This is the part supplied by the classical H\"ormander calculus and its semiclassical formulations \citep{hormander1985analysis3,hormander1985analysis4,zworski2012semiclassical,alazardmicrolocal}. The rest of the proof only uses this parametrix through the generated phase, amplitude symbol bounds, and the absence of competing stationary branches.

\paragraph{Step 2: test against packets.}
Insert $\varphi_\eta^h$ and $\varphi_\gamma^h$ into the kernel representation. The resulting phase is the half-wave phase plus the packet phases centered at $z_\eta$ and $z_\gamma$.

\paragraph{Step 3: phase-gradient lower bound outside the canonical tube.}
If $z_\gamma$ lies outside a fixed canonical tube around $\kappa_t(z_\eta)$, then the packet-localized total phase has gradient bounded from below by
\[
|\nabla\Psi_{t,\gamma,\eta}|
\ge
c_T h^{-1/2}\,\mathrm{dist}\bigl(z_\gamma,\kappa_t(z_\eta)\bigr)-C_T\delta_h
\]
on the support of the localized amplitude, for constants depending on the compact phase-space window and the time horizon $T$. After enlarging the canonical tube by a fixed multiple of the normalized covering defect $\delta_h$, the last term is absorbed and the lower bound becomes
\[
|\nabla\Psi_{t,\gamma,\eta}|
\ge
c'_T h^{-1/2}\,\mathrm{dist}\bigl(z_\gamma,\kappa_t(z_\eta)\bigr).
\]
This is exactly where the no-caustic hypothesis is used: the canonical graph is smooth, locally single-valued, and the phase has no competing nearby stationary point away from the propagated packet center.

\paragraph{Step 4: repeated integration by parts.}
Applying the non-stationary phase operator
\[
L
=
\frac{\nabla\Psi_{t,\gamma,\eta}\cdot\nabla}{i|\nabla\Psi_{t,\gamma,\eta}|^2}
\]
repeatedly transfers derivatives to the localized amplitude. Each application gains one factor of
\[
\Bigl(1+h^{-1/2}\,\mathrm{dist}\bigl(z_\gamma,\kappa_t(z_\eta)\bigr)\Bigr)^{-1},
\]
and the analytic packet envelope keeps the differentiated amplitude uniformly summable. After $N$ iterations one obtains~\eqref{eq:offgraph}.
\end{proof}

Theorem~\ref{thm:wave-offgraph} is a classical manuscript-level theorem and the first continuous spine of the paper.

\begin{proposition}[Restricted non-stationary phase bridge]
\label{prop:restricted-nsp-bridge}
Fix a packet-localized oscillatory representation of a canonical-graph propagator on a short time window. Assume only the following restricted inputs:
\begin{enumerate}
\item a single-graph oscillatory parametrix on the chosen frequency band;
\item uniform derivative bounds for the localized amplitude after packet testing;
\item a phase-gradient lower bound outside a raw canonical tube of radius $O(h^{1/2}(1+\delta_h))$.
\end{enumerate}
Then repeated integration by parts yields the off-graph decay law~\eqref{eq:offgraph}. In particular, the manuscript does not need the full H\"ormander calculus at the point where the packet theorem starts. It needs only this restricted bridge from canonical-graph geometry to packet-matrix decay.
\end{proposition}

\begin{proof}
The proposition is the rescaled packet form of Steps~3--4 in the proof of Theorem~\ref{thm:wave-offgraph}.  A full statement and proof are given in Proposition~\ref{prop:restricted-nsp-appendix}.  In brief, once the phase-gradient lower bound is available on the packet-localized support, the operator
\[
L=\frac{\nabla\Psi\cdot\nabla}{i|\nabla\Psi|^2}
\]
can be iterated without encountering a stationary point.  After rescaling packet variables by $h^{1/2}$, the large parameter is
\[
D_{\gamma\eta}=1+h^{-1/2}\mathrm{dist}(z_\gamma,\kappa(z_\eta)).
\]
The differentiated coefficients of the adjoint vector field are $O(D_{\gamma\eta}^{-1})$, while the rescaled packet amplitude has uniformly bounded $C^N$-$L^1$ seminorms.  Iterating $N$ times therefore gives a factor $D_{\gamma\eta}^{-N}$, which is exactly~\eqref{eq:offgraph}.
\end{proof}

\subsection{Spine B: abstract canonical-graph packet sparsity}

The second spine abstracts away from the wave equation and isolates the exact packet argument that converts off-graph decay into bounded-stencil sparsity.

\begin{theorem}[Abstract canonical-graph packet sparsity]
\label{thm:packet-sparsity}
Let $F_h$ be a packet matrix associated with a canonical graph $\kappa$, and suppose that for every $N\in\N$ there is $C_N>0$ such that
\[
\bigl|\langle F_h\varphi_\eta^h,\varphi_\gamma^h\rangle\bigr|
\le
C_N \Bigl(1+h^{-1/2}\,\mathrm{dist}\bigl(z_\gamma,\kappa(z_\eta)\bigr)\Bigr)^{-N}.
\]
Assume in addition that there exist constants $q>0$ and $C_{\mathrm{cnt}}>0$ such that for every source packet index $\eta$ and every shell radius $R\ge 1$,
\begin{equation}
\label{eq:packet-count}
\#\Bigl\{\gamma\in\Gamma_h :
\mathrm{dist}\bigl(z_\gamma,\kappa(z_\eta)\bigr)\le R h^{1/2}\Bigr\}
\le
C_{\mathrm{cnt}} R^q.
\end{equation}
Let $F_h^{(K)}$ be the packet matrix obtained by keeping, for each source packet, only the $K$ output packets nearest to $\kappa(z_\eta)$.  Assume the same shell-counting estimate for the inverse graph, equivalently the retained matrix and its adjoint have the same tail envelope on the compact no-caustic window. Then for every $N>q+1$ there is a constant $C_{N,q}$ such that
\[
\|F_h-F_h^{(K)}\|_{\ell^2\to\ell^2}
\le
C_{N,q}\, K^{1-N/q}.
\]
\end{theorem}

\begin{proof}[Proof sketch]
Decompose the packet indices into dyadic shells around the propagated center $\kappa(z_\eta)$. On shell $\ell$, the counting hypothesis gives $O(2^{q\ell})$ packets and the off-graph decay gives $O(2^{-N\ell})$ per entry. Hence the column shell mass is $O(2^{(q-N)\ell})$. The inverse-graph hypothesis gives the same row shell mass. Summing over shells above level $L$ gives row and column tails of size $O(2^{(q-N)L})$. Since the number of packets retained up to shell $L$ is $O(2^{qL})$, one has $K\asymp 2^{qL}$; Schur's test gives the displayed $\ell^2$ bound.
\end{proof}

Theorem~\ref{thm:packet-sparsity} is the manuscript's abstract FIO theorem. It is independent of the wave equation, and Theorem~\ref{thm:wave-offgraph} shows that short-time half-wave propagation satisfies its hardest hypothesis. The corresponding Lean files formalize the shell/tube/truncation layer that starts once the off-graph decay and shell counting laws are available.

\begin{lemma}[Row-column bounded-stencil transfer]
\label{lem:row-column-stencil}
Assume, in addition to the shell-counting hypothesis above, that $\kappa$ is bi-Lipschitz on the compact packet window and that the packet cover has bounded overlap at scale $h^{1/2}$.  If a source-wise truncation keeps at most $K$ output packets in a canonical tube around $\kappa(z_\eta)$, then there is a constant $C_{\mathrm{deg}}$, depending only on the bi-Lipschitz and overlap constants, such that every output packet receives contributions from at most $C_{\mathrm{deg}}K$ retained source packets.  Conversely, an output-wise canonical tube stencil has at most $C_{\mathrm{deg}}K$ retained outputs per source.
\end{lemma}

\begin{proof}[Proof sketch]
Fix an output packet $\gamma$.  If $\gamma$ is retained by a source packet $\eta$, then $z_\gamma$ lies in a tube of radius $R_Kh^{1/2}$ around $\kappa(z_\eta)$, where $R_K$ is the shell radius corresponding to $K$ retained packets.  Since $\kappa^{-1}$ is Lipschitz on the compact window, $z_\eta$ lies in a tube of radius $C R_Kh^{1/2}$ around $\kappa^{-1}(z_\gamma)$.  The bounded-overlap shell-counting law bounds the number of such source packets by $C_{\mathrm{deg}}R_K^q$.  Because the retained source-wise tube has $K\asymp R_K^q$ packets, this is $O(K)$.  The reverse implication is identical with $\kappa$ and $\kappa^{-1}$ interchanged.
\end{proof}

Lemma~\ref{lem:row-column-stencil} resolves the apparent direction mismatch between the abstract sparsity theorem and the executable layer.  The sparsity proof is most naturally columnwise, retaining output packets for each source packet.  The \MiNO\ implementation is outputwise, aggregating source packets for each output packet.  On compact no-caustic windows these two descriptions differ only by the bounded degree constant above, which is absorbed into the stencil budget and constants below.

\subsection{Continuous transported-symbol reduction with bounded canonical stencil}

We now return to the actual operator class used by \MiNO. The previous theorem says that a canonical-graph packet matrix can be truncated to a bounded stencil with explicit error. The next result combines that with stationary phase near the graph.

\begin{theorem}[Continuous transported-symbol reduction]
\label{thm:continuous-reduction}
Under the assumptions of Theorem~\ref{thm:wave-offgraph}, fix an integer $M\ge 1$ and a stencil budget $K\in\N$. Then the packet-space realization
\[
K_h(t):=W_h U_+(t) W_h^\ast
\]
admits a decomposition
\[
K_h(t)={\transportref_h}^{(K)}(t)+R_h^{(K,M)}(t),
\]
where
\[
({\transportref_h}^{(K)}(t)a)_\gamma
=
\sum_{\eta\in\Pi^\star_{h,K}(t,\gamma)}
s_{h,\gamma\eta}^\star(t)\,a_\eta
\]
has an output-wise canonical stencil of size at most $C_{\mathrm{deg}}K$ after the row-column transfer of Lemma~\ref{lem:row-column-stencil} (renamed as $K$ below), with each retained packet center lying in a raw tube of radius $O(h^{1/2}(1+\delta_h))$ around the canonical preimage, and where the residual satisfies
\[
\|R_h^{(K,M)}(t)\|_{\ell^2\to\ell^2}
\le
C_{T,N,q}K^{1-N/q}
+E_{\mathrm{frame}}(h)+E_{\mathrm{asymp}}(h,M)+C_{\mathrm{cov}}\delta_h
\]
for every $N>q+1$.
\end{theorem}

\begin{proof}[Proof sketch]
Apply Theorem~\ref{thm:packet-sparsity} to discard all but the $K$ nearest packets to the canonical graph in the source-to-output direction. Lemma~\ref{lem:row-column-stencil} converts this into the output-to-source stencil used by \MiNO, changing only constants and the effective value of $K$. This contributes the $K^{1-N/q}$ tail. On the retained tube, stationary phase around the propagated packet center identifies the leading amplitude samples
\[
s_{h,\gamma\eta}^\star(t)
=
a_t(z_\eta,z_\gamma)+O(h^{1/2}(1+\delta_h)),
\]
while bounded overlap of the packet frame and Schur's test convert the local coefficient error into an $\ell^2$ operator remainder.  We separate this finite-to-asymptotic term into the frame defect $E_{\mathrm{frame}}(h)$, the retained stationary-phase defect $E_{\mathrm{asymp}}(h,M)$, and the normalized covering defect $C_{\mathrm{cov}}\delta_h$. Summing the truncation and local approximation terms gives the stated bound.
\end{proof}

Theorem~\ref{thm:continuous-reduction} is the continuous theorem that the paper now wants to stand behind. It is stronger than the old transported-symbol statement because it identifies the correct \emph{local rank} of the reference operator in the packet energy norm: not dense, not necessarily one, but bounded by a packet sparsity budget $K$.

Theorem~\ref{thm:continuous-reduction} is also the paper's restricted packet almost-diagonalization bridge. The off-graph part comes from Proposition~\ref{prop:restricted-nsp-bridge} plus shell counting, and the on-graph part comes from a local stationary-phase sample of the retained amplitude. This is the exact point where the continuous argument becomes a theorem-grade packet reduction rather than a broad appeal to ``standard FIO theory.'' The remaining classical input is therefore narrow and inspectable.

\begin{corollary}[Continuous-to-discrete bounded-stencil \MiNO\ corollary]
\label{cor:continuous-to-discrete}
Under the assumptions of Theorem~\ref{thm:continuous-reduction}, let ${\widehat{\transport}_h}^{(K)}(t)$ be a theorem-grade \MiNO\ block with the same stencil budget $K$. This corollary is the finite landing statement: it assumes coefficient-level budgets after a retained stencil has already been matched. The main wave theorem below replaces this finite transport hypothesis by a geometric canonical-map error and the induced packet-kernel mismatch. Suppose that for every retained pair $(\gamma,\eta)$ one has transport and symbol budgets
\[
|a_{\hat\pi_{\gamma\eta}}-a_{\pi^\star_{\gamma\eta}}|\le \eps_{\mathrm{tr}},
\qquad
|\hat s_{h,\gamma\eta}-s_{h,\gamma\eta}^\star|\le \eps_{\mathrm{sym}},
\]
and residual envelope $|\rho_{h,\gamma}(a)|\le \eps_{\mathrm{res}}$, together with coefficient/symbol envelopes $|a_\eta|\le B_a$ and $|s_{h,\gamma\eta}^\star|\le B_s$. Then
\begin{align*}
\|{\widehat{\transport}_h}^{(K)}(t)a-K_h(t)a\|_{\ell^1}
\le{}&
nK(B_s\eps_{\mathrm{tr}}+\eps_{\mathrm{sym}}B_a)
+n\eps_{\mathrm{res}}\\
&+
\Bigl(C_{T,N,q}K^{1-N/q}
+E_{\mathrm{frame}}(h)+E_{\mathrm{asymp}}(h,M)+C_{\mathrm{cov}}\delta_h\Bigr)\|a\|_{\ell^1}.
\end{align*}
\end{corollary}

\begin{proof}
The retained stencil part is bounded by applying the finite propagation bookkeeping packetwise across at most $K$ retained interactions per output packet. The discarded shells and packet-covering defect are then absorbed into the residual term from Theorem~\ref{thm:continuous-reduction}.
\end{proof}

\subsection{Network-realizable short-time wave approximation}
\label{sec:network-realizable-wave}

Corollary~\ref{cor:continuous-to-discrete} assumes small learned transport and symbol budgets.  The next theorem closes that gap at the level of approximation theory: under the same no-caustic wave assumptions, a theorem-grade \MiNO-Core block with sufficiently large metadata and symbol networks realizes the transported-symbol approximation with an explicit rate.

Let $\Omega_T\subset T^\ast\R^d$ be the compact packet window swept out by the no-caustic half-wave flow for $|t|\le T$, and let $m_\Omega=\dim(\Omega_T)$.  We use the following standard compact-window neural approximation input: if $f\in C^r(\Omega_T;\R^p)$ with $\|f\|_{C^r}\le A_r$, then the branch network class with $P$ trainable parameters contains a map $\hat f_P$ such that
\begin{equation}
\label{eq:branch-approx-rate}
\|f-\hat f_P\|_{L^\infty(\Omega_T)}
\le
C_{\mathrm{nn}}A_r P^{-r/m_\Omega},
\end{equation}
up to the usual architecture-dependent logarithmic factors.  For ReLU networks this is the standard compact smooth-function approximation rate \citep{yarotsky2017relu}; other smooth activations may be substituted without changing the structure of the theorem.

\begin{proposition}[Hamiltonian metadata branch consistency]
\label{prop:hamiltonian-metadata-consistency}
Let $H\in C^{r+2}(\Omega_T)$ generate the canonical flow $\kappa_t$ by
\[
\dot z = X_H(z)=J\nabla H(z),
\qquad
z=(x,\xi),
\]
on the compact no-caustic window $\Omega_T$, and suppose $\|H\|_{C^{r+2}}\le A_{r+2}$.  Assume a scalar Hamiltonian branch realizes $H_P$ with
\[
\|H-H_P\|_{C^2(\Omega_T)}
\le
C_H A_{r+2}P_H^{-(r-2)/m_\Omega},
\qquad r>2.
\]
Let $\widehat\kappa_{t,P,\Delta}$ be the $m$-step explicit Hamiltonian-vector-field update
\[
z_{j+1}=z_j+\Delta\,J\nabla H_P(z_j),
\qquad
\Delta=t/m,
\]
projected back to $\Omega_T$ when needed.  Then, for $|t|\le T$,
\[
\sup_{z\in\Omega_T}
d\bigl(\widehat\kappa_{t,P,\Delta}(z),\kappa_t(z)\bigr)
\le
C_T\Bigl(P_H^{-(r-2)/m_\Omega}+\Delta\Bigr).
\]
Moreover one explicit step has canonical-defect bound
\[
\sup_{z\in\Omega_T}
\bigl\|D\Psi_{P,\Delta}(z)^\top\Omega_0D\Psi_{P,\Delta}(z)-\Omega_0\bigr\|
\le
C_T\Delta^2\|H_P\|_{C^2(\Omega_T)}^2,
\]
where $\Psi_{P,\Delta}(z)=z+\Delta J\nabla H_P(z)$ and $\Omega_0$ is the standard symplectic matrix.

If, in addition, $H_P$ is represented in separable form
\[
H_P(q,p)=V_P(q)+T_P(p),
\]
and the metadata branch uses the St\"ormer--Verlet step
\[
p_{1/2}=p_0-\frac{\Delta}{2}\nabla_qV_P(q_0),\qquad
q_1=q_0+\Delta\nabla_pT_P(p_{1/2}),\qquad
p_1=p_{1/2}-\frac{\Delta}{2}\nabla_qV_P(q_1),
\]
then each step is symplectic in the metadata variables.  The canonical-defect term vanishes at the map level, up to floating-point and autograd implementation error, and the discretization term in the flow error is the standard second-order $O(\Delta^2)$ term for smooth separable Hamiltonians.
\end{proposition}

\begin{proof}
The vector-field error is
\[
\|X_H-X_{H_P}\|_{C^1}
\le
C\|H-H_P\|_{C^2}.
\]
Standard Gronwall stability for ODE flows on the compact window gives an $O(\|X_H-X_{H_P}\|_{C^0})$ flow perturbation over $|t|\le T$.  The explicit Euler discretization for the $C^1$ vector field $X_{H_P}$ contributes the usual first-order global error $O(\Delta)$, with a constant depending on the $C^1$ and Lipschitz bounds of $X_{H_P}$.  Combining the two estimates gives the displayed flow error.  For the canonical defect, write
\[
D\Psi_{P,\Delta}=I+\Delta\,DX_{H_P}.
\]
Since $DX_{H_P}=J\nabla^2H_P$ is an infinitesimal Hamiltonian matrix, the first-order term in
\[
(I+\Delta DX_{H_P})^\top\Omega_0(I+\Delta DX_{H_P})-\Omega_0
\]
cancels.  The remaining term is quadratic in $\Delta$ and bounded by $C\Delta^2\|DX_{H_P}\|^2$, hence by $C_T\Delta^2\|H_P\|_{C^2}^2$ on the compact window.

For the separable branch, the St\"ormer--Verlet map is a composition of three exact symplectic shears: a potential half-step, a kinetic full-step, and another potential half-step. A composition of symplectic maps is symplectic, so the discrete canonical defect is zero at the exact map level. The trajectory approximation is the standard second-order consistency estimate for smooth separable Hamiltonian flows on a compact window.
\end{proof}

Proposition~\ref{prop:hamiltonian-metadata-consistency} records the scope of the Hamiltonian metadata branch.  It justifies \texttt{hamiltonian\_euler} as a cheap Hamiltonian vector-field bias and \texttt{hamiltonian\_verlet} as the theorem-facing separable-Hamiltonian transport branch.  The latter is symplectic by construction in the packet metadata variables under the stated separability and fixed-feature assumptions.  The wave theorem below uses the broader compact-window transport approximation hypothesis; the Verlet branch is a clean structural specialization of that hypothesis.

\begin{theorem}[Network-realizable short-time wave \MiNO-Core approximation]
\label{thm:network-realizable-wave-mino}
Assume the hypotheses of Theorems~\ref{thm:wave-offgraph}, \ref{thm:packet-sparsity}, and \ref{thm:continuous-reduction}.  Assume that the canonical map $\kappa_t:\Omega_T\to\Omega_T$ has $C^r$ norm at most $A_r$ on the compact packet window, with $r\ge 1$.  Assume also that the retained local packet kernel has the form
\[
s_t=s_t(z,\nu;h),
\qquad
\nu=h^{-1/2}\bigl(z_\gamma-\kappa_t(z_\eta)\bigr),
\]
with $C^r$ norm at most $A_r$ on a compact tube-coordinate window of dimension $m_\Omega+q_{\mathrm{loc}}$.  The source-only multiplier case is recovered by taking $q_{\mathrm{loc}}=0$.  Let the metadata and symbol branches have budgets $P_{\mathrm{tr}}$ and $P_{\mathrm{sym}}$, both satisfying~\eqref{eq:branch-approx-rate} on their respective input dimensions.  Then there exists a theorem-grade bounded-stencil \MiNO-Core block $\mathcal M_{h,K,P}^{\mathrm{core}}(t)$ such that, for every $N>q+1$ and every retained asymptotic order $M\ge 1$,
\begin{align}
\label{eq:network-wave-rate}
\|W_hU_+(t)W_h^\ast-\mathcal M_{h,K,P}^{\mathrm{core}}(t)\|_{\ell^2\to\ell^2}
\le{}&
C_{\mathrm{tail}}K^{1-N/q}
+E_{\mathrm{frame}}(h)
+E_{\mathrm{asymp}}(h,M)
+C_{\mathrm{cov}}\delta_h \notag\\
&+C_{\mathrm{tr}}K\bigl(h^{-1/2}P_{\mathrm{tr}}^{-r/m_\Omega}+\delta_h\bigr)
 +C_{\mathrm{sym}}K P_{\mathrm{sym}}^{-r/(m_\Omega+q_{\mathrm{loc}})}\notag\\
&+C_{\mathrm{res}}\eps_{\mathrm{res}} .
\end{align}
The constants depend on the time window, packet frame bounds, shell-counting geometry, compact-window $C^r$ norms, local tube dimension, and the branch architecture class, but not on $h$, $K$, $P_{\mathrm{tr}}$, or $P_{\mathrm{sym}}$.  The normalized covering defect $\delta_h$ is an explicit packetization term: it vanishes only for an oversampled packet cloud with normalized covering radius tending to zero.  For a critically sampled frame, $\delta_h$ should be read as a fixed finite-frame defect rather than as an asymptotically vanishing quantity.
\end{theorem}

\begin{proof}[Proof sketch]
The proof combines five ingredients.

\paragraph{Step 1: reduce the PDE operator to a sparse packet matrix.}
Theorem~\ref{thm:wave-offgraph} gives off-canonical-tube decay for the half-wave propagator.  Theorem~\ref{thm:packet-sparsity} turns that decay into the $K$-term packet tail $C_{\mathrm{tail}}K^{1-N/q}$.  Theorem~\ref{thm:continuous-reduction} then identifies the retained part with a bounded-stencil transported-symbol packet law, up to the frame, asymptotic, and covering defects $E_{\mathrm{frame}}(h)+E_{\mathrm{asymp}}(h,M)+C_{\mathrm{cov}}\delta_h$.

\paragraph{Step 2: approximate the canonical map.}
Apply~\eqref{eq:branch-approx-rate} to $\kappa_t$.  The metadata branch can therefore realize a map $\widehat\kappa_t$ with
\[
\|\widehat\kappa_t-\kappa_t\|_{L^\infty(\Omega_T)}
\le
C_{\mathrm{nn}}A_rP_{\mathrm{tr}}^{-r/m_\Omega}.
\]
Because packet tubes have natural radius $h^{1/2}$ in the semiclassical scaling, this geometric error enters the packet coefficient law as
\[
\eps_{\mathrm{tr}}
\le
C_\kappa\bigl(h^{-1/2}P_{\mathrm{tr}}^{-r/m_\Omega}+\delta_h\bigr).
\]
The $h^{-1/2}$ factor is the cost of measuring canonical-map error in packet-radius units.

\paragraph{Step 3: approximate the local packet kernel.}
Apply~\eqref{eq:branch-approx-rate} to the retained local packet kernel $s_t(z,\nu;h)$ on the compact source--tube coordinate window.  The symbol branch can realize $\widehat s_t$ with
\[
\eps_{\mathrm{sym}}
\le
C_s P_{\mathrm{sym}}^{-r/(m_\Omega+q_{\mathrm{loc}})}.
\]

\paragraph{Step 4: insert the operator-level bounded-stencil bridge.}
The retained packet law has at most $K$ interactions per output packet after Lemma~\ref{lem:row-column-stencil}.  The geometric canonical-map error first gives a packet-kernel transport mismatch of size $O(h^{-1/2}P_{\mathrm{tr}}^{-r/m_\Omega}+\delta_h)$ on the retained tube.  Theorem~\ref{thm:stencil-propagation} then converts that kernel mismatch, together with the symbol approximation error, into the learned-budget contribution
\[
C_{\mathrm{tr}}K\eps_{\mathrm{tr}}
+C_{\mathrm{sym}}K\eps_{\mathrm{sym}}
+C_{\mathrm{res}}\eps_{\mathrm{res}}.
\]
Equivalently, Theorem~\ref{thm:schur-packet-matrix} gives the same bounded-stencil contribution directly in $\ell^2\to\ell^2$ operator norm after row and column stencil bounds are imposed.

\paragraph{Step 5: combine defects.}
Summing the tail, frame, asymptotic, covering, transport, local-kernel, and residual terms gives~\eqref{eq:network-wave-rate}.
\end{proof}

Theorem~\ref{thm:network-realizable-wave-mino} is the rate-realization input for the main Microlocal Propagation Theorem stated in Section~\ref{sec:pde}.  Its scope is compact-window, no-caustic, short-time hyperbolic propagation, with classical off-graph and stationary-phase inputs and a finite machine-checked landing layer.  Within that scope, the theorem proves that the \MiNO-Core primitive approximates a concrete PDE operator family in packet energy norm, with an explicit rate in packet scale, stencil budget, covering radius, local-kernel dimension, and neural-network size.

\begin{corollary}[Sobolev field-level MiNO-Core approximation]
\label{cor:sobolev-field-mpt}
Assume the hypotheses of Theorem~\ref{thm:network-realizable-wave-mino}.  In addition, assume that the packet analysis and synthesis maps are Sobolev-stable on the compact phase window: for every $s$ in a fixed compact interval $I\subset\mathbb R$ there are constants $A_{W,s},B_{W,s}$, independent of $h$ and the packet count, such that the weighted packet norms induced by the window satisfy
\[
\|W_h v\|_{\ell^2_s}\le A_{W,s}\|v\|_{H^s},
\qquad
\|W_h^\ast b\|_{H^s}\le B_{W,s}\|b\|_{\ell^2_s},
\]
and the retained short-time canonical graph changes the corresponding weights by at most a constant $C_{\kappa,s}$ on $\Omega_T$.  Then, for microlocalized inputs and for every $s\in I$,
\[
\|W_h^\ast\mathcal M_{h,K,P}^{\mathrm{core}}(t)W_h-U_+(t)\|_{H^s\to H^s}
\le
C_s\bigl(
E_{h,K,P}
+E^{\mathrm{in}}_{\mathrm{proj},s}
+E^{\mathrm{out}}_{\mathrm{proj},s}
\bigr),
\]
where $E_{h,K,P}$ denotes the right-hand side of Theorem~\ref{thm:network-realizable-wave-mino}, and
\[
E^{\mathrm{in}}_{\mathrm{proj},s}
=
\|U_+(t)(I-\mathcal P_h)\|_{H^s\to H^s},
\qquad
E^{\mathrm{out}}_{\mathrm{proj},s}
=
\|(I-\mathcal P_h)U_+(t)\mathcal P_h\|_{H^s\to H^s}.
\]
The same statement with the executable transported-synthesis, carrier, and admissible landing defects added holds for the carrier-bound field realization.
\end{corollary}

\begin{proof}
Apply Theorem~\ref{thm:network-realizable-wave-mino} in the weighted packet norm $\ell^2_s$.  The compact-window canonical graph preserves Sobolev weights up to $C_{\kappa,s}$, and the retained symbol is order zero for the half-wave branch, so the same packet error envelope $E_{h,K,P}$ applies with a Sobolev-dependent constant.  Sobolev-stable synthesis converts the packet error to $H^s$, while analysis stability converts the input $H^s$ norm to the weighted packet norm.  The two displayed projection defects account for replacing the true field by its packet projection before and after propagation.  The executable transported-synthesis, carrier, and finite-landing terms enter by the same triangle inequality as in the $L^2$ field-level statement.
\end{proof}

\subsection{Branch-factored oscillatory extension}
\label{sec:branch-factored-extension}

The preceding theorem treats a single canonical graph.  This is the cleanest short-time wave setting, but it should not be confused with the full range of oscillatory PDE operators.  Local outgoing Helmholtz windows, reflected waves, and scattering parametrices may be better modeled by a finite sum of canonical-graph FIOs, while elliptic and parabolic regimes remain identity-canonical or smoothing components.  The architectural extension is therefore not Helmholtz-specific: it is a finite-branch packet realization of a local FIO sum.

Let an oscillatory target on a compact packet window admit the microlocal decomposition
\[
T_h
=
\sum_{\ell=1}^{L} T_{\ell,h}
+A_h+S_h,
\]
where each $T_{\ell,h}$ is a canonical-graph FIO with graph $\kappa_\ell$, retained symbol $s_\ell$, and local packet tube, $A_h$ is an identity-canonical pseudodifferential component, and $S_h$ is smoothing or outside the retained window.  A branch-factored \MiNO\ realizes
\[
\mathcal M_{\theta,h}^{(L)}u
=
W_h^\ast
\left[
\sum_{\ell=1}^{L}
G_{\theta,\ell}(z)\,
\mathcal T_{\hat\kappa_{\theta,\ell},\hat s_{\theta,\ell},\hat\phi_{\theta,\ell}}(W_hu)
+\mathcal A_\theta(W_hu)
+R_\theta(W_hu)
\right],
\]
with token-wise router $\sum_{\ell}G_{\theta,\ell}=1$.  The optional phase factor $\hat\phi_{\theta,\ell}$ is included only in regimes where the packet landing error is dominated by unresolved high-frequency phase.  It does not change the single-branch wave theorem; it records the sharper implementation requirement that phase errors scale like $h^{-1}$ in oscillatory output.

\begin{lemma}[Phase-factor landing defect]
\label{lem:phase-factor-defect}
On a local output window, suppose a branch is represented as
\[
T_{\ell,h}u=e^{i\phi_\ell/h}B_{\ell,h}u,
\qquad
\widehat T_{\ell,h}u=e^{i\hat\phi_\ell/h}\widehat B_{\ell,h}u,
\]
with $\|B_{\ell,h}\|_{L^2\to L^2}\le C_B$ and $\|u\|_{L^2}\le 1$.  Then
\[
\|(T_{\ell,h}-\widehat T_{\ell,h})u\|_{L^2}
\le
\|(B_{\ell,h}-\widehat B_{\ell,h})u\|_{L^2}
+C_B h^{-1}\|\phi_\ell-\hat\phi_\ell\|_{L^\infty}.
\]
The same estimate holds in packet norm after multiplying the right-hand side by the packet frame synthesis constant.
\end{lemma}

\begin{proof}
Use
$|e^{ia/h}-e^{ib/h}|\le h^{-1}|a-b|$
and the triangle inequality.  The packet statement follows by applying the same synthesis bound used in Corollary~\ref{cor:sobolev-field-mpt}.
\end{proof}

\begin{theorem}[Branch-factored oscillatory \MiNO\ envelope]
\label{thm:branch-factored-oscillatory-mino}
Assume that $T_h$ admits the finite decomposition above, that the canonical tubes of the retained FIO branches have finite overlap, and that the ideal router $G^\star_\ell$ assigns retained packets to the corresponding branch up to router defect $E_{\mathrm{route}}$.  For branch $\ell$, let
$E_{\mathrm{tr},\ell}$, $E_{\mathrm{sym},\ell}$, and $E_{\mathrm{phase},\ell}$ denote its transport, symbol, and optional phase-factor defects; let $E_{\mathrm{tsyn},\ell}$, $E_{\mathrm{car},\ell}$, and $E_{\mathrm{land},\ell}$ denote the executable transported-synthesis, carrier, and finite-landing defects; and let $E_{\mathrm{cross}}$ denote packet cross-talk between retained branch tubes.  Then the branch-factored carrier-bound realization satisfies
\[
\begin{aligned}
\|\mathcal M_{\theta,h}^{(L)}-T_h\|_{\ell^2\to\ell^2}
&\le
\sum_{\ell=1}^{L}
\bigl(
E_{\mathrm{tr},\ell}
+E_{\mathrm{sym},\ell}
+E_{\mathrm{phase},\ell}
+E_{\mathrm{tsyn},\ell}
+E_{\mathrm{car},\ell}
+E_{\mathrm{land},\ell}
\bigr)\\
&\quad
+E_{\mathrm{route}}
+E_{\mathrm{cross}}
+E_{\Psi}
+E_{\mathrm{res}}
+E_{\mathrm{tail}}
+E_{\mathrm{frame}}+E_{\mathrm{asymp}}+C_{\mathrm{cov}}\delta_h .
\end{aligned}
\]
Here $E_{\Psi}$ is the identity-canonical pseudodifferential approximation defect for $A_h$.  The corresponding field-level $L^2$ and Sobolev estimates add the same projection constants as Corollary~\ref{cor:sobolev-field-mpt}.
\end{theorem}

\begin{proof}
Apply the single-branch packet reduction to each retained canonical tube and sum the per-branch transport, symbol, phase, and executable landing defects.  Replacing the ideal branch assignment by the learned router gives $E_{\mathrm{route}}$.  Finite overlap and almost-orthogonality of the retained tubes control the non-diagonal interactions by $E_{\mathrm{cross}}$.  The identity-canonical component is bounded by the symbol-order approximation argument used for the $\Psi$DO branch, producing $E_{\Psi}$.  Smoothing and out-of-window terms are absorbed into $E_{\mathrm{res}}$, $E_{\mathrm{tail}}$, and the frame/asymptotic/covering defects.  The field-level statement is the same analysis--synthesis triangle inequality as before.
\end{proof}

Theorem~\ref{thm:branch-factored-oscillatory-mino} is the correct home for high-frequency stress tests.  It shows that branch count is only one term in the budget: improvement also requires small per-branch transport and symbol errors, a router that separates canonical tubes, limited cross-talk, and phase/landing accuracy that tightens with frequency.  A flat branched Helmholtz table therefore identifies a finite-budget failure of these hypotheses while leaving the single-graph wave theorem intact.

\paragraph{Finite packet Egorov diagnostic.}
The next microlocal calculus question is whether $\mathcal M_{h,K,P}^{\mathrm{core}}$ not only approximates $U_+(t)$ but also conjugates pseudodifferential probes by the learned canonical map.  For a frozen theorem-grade linearized Core, let $A_{j,h}^{\mathrm{pkt}}=W_h\Op_h(a_j)W_h^\ast$ be a finite family of order-zero packet probes, let $A_{j,h}^{\widehat\kappa}=W_h\Op_h(a_j\circ\widehat\kappa_\theta)W_h^\ast$, and let $S_{\theta,h}$ be the retained diagonal or local symbol action of the learned symbol branch.  Define
\[
\rho_{\mathrm{Eg},h}^{\mathrm{pkt}}(j)
=
\frac{
\bigl\|
(\mathcal M_{\theta,h}^{\mathrm{core}})^\ast
A_{j,h}^{\mathrm{pkt}}
\mathcal M_{\theta,h}^{\mathrm{core}}
-
S_{\theta,h}^\ast A_{j,h}^{\widehat\kappa}S_{\theta,h}
\bigr\|_{\ell^2\to\ell^2}
}{
\|A_{j,h}^{\mathrm{pkt}}\|_{\ell^2\to\ell^2}+\epsilon
}.
\]
For a genuinely nonlinear executable model, the corresponding diagnostic must be applied to the local linearization $D\mathcal M_\theta[u]$ rather than to $\mathcal M_\theta$ itself.  Thus the planned Egorov metric is
\[
\rho_{\mathrm{Eg},h}^{\mathrm{Jac}}(u,j)
=
\frac{
\bigl\|
D\mathcal M_\theta[u]^\ast A_{j,h}^{\mathrm{pkt}}D\mathcal M_\theta[u]
-
S_{\theta,h}[u]^\ast A_{j,h}^{\widehat\kappa[u]}S_{\theta,h}[u]
\bigr\|_{\ell^2\to\ell^2}
}{
\|A_{j,h}^{\mathrm{pkt}}\|_{\ell^2\to\ell^2}+\epsilon
}.
\]
This diagnostic records the finite object for branch-coupled Egorov measurements.  The main closure of this paper remains packet propagation and finite wavefront proxies.

\begin{corollary}[Packet-threshold wavefront localization]
\label{cor:packet-wavefront-localization}
Fix a packet threshold $\tau>0$ and define the discrete $h$-wavefront proxy
\[
\WF^\tau_h(a)=\{z_\eta:\ |a_\eta|\ge \tau\}.
\]
Let $a=W_hu$ satisfy $\|a\|_{\ell^1}\le A$ and let $S=\WF^\tau_h(a)$.  Under the hypotheses of Theorem~\ref{thm:network-realizable-wave-mino}, the coefficients of
\[
\mathcal M_{h,K,P}^{\mathrm{core}}(t)a
\]
outside the raw phase-space neighborhood of $\kappa_t(S)$ obey
\begin{align*}
\sum_{z_\gamma\notin (\kappa_t(S))_{h^{1/2}(R_K+C\epsilon_{\kappa,h}+C\delta_h)}}
\bigl|(\mathcal M_{h,K,P}^{\mathrm{core}}(t)a)_\gamma\bigr|
\le{}&
C A K^{1-N/q}
C A\bigl(E_{\mathrm{frame}}(h)+E_{\mathrm{asymp}}(h,M)\bigr)\\
&+C A C_{\mathrm{cov}}\delta_h
+A E_{\mathrm{net}}
+C\tau ,
\end{align*}
where
\[
\epsilon_{\kappa,h}=h^{-1/2}\sup_{z\in\Omega_T}d(\widehat\kappa_\theta(z),\kappa_t(z)),
\qquad
E_{\mathrm{net}}
=
K\bigl(h^{-1/2}P_{\mathrm{tr}}^{-r/m_\Omega}+\delta_h\bigr)
+K P_{\mathrm{sym}}^{-r/(m_\Omega+q_{\mathrm{loc}})}
+\epsilon_{\mathrm{res}} .
\]
\end{corollary}

\begin{proof}
Split the input coefficients into the threshold-active part supported on $S$ and the remainder of $\ell^1$ mass at most $C\tau$ after restricting to the finite active packet window.  The active part is transported into the canonical tube around $\kappa_t(S)$ by the retained-stencil construction; the raw tube radius is enlarged by $h^{1/2}\epsilon_{\kappa,h}$ for learned canonical-map error, by $h^{1/2}\delta_h$ for packet covering defect, and by the retained shell radius $R_Kh^{1/2}$.  Contributions outside that enlarged tube are exactly the discarded off-graph shells, packetization error, and learned-budget error from Theorem~\ref{thm:network-realizable-wave-mino}.  Summing these terms gives the displayed estimate.
\end{proof}

Classically, $(x,\xi)\notin \WF_h(u)$ means that some semiclassical cutoff supported near $(x,\xi)$ annihilates $u$ up to $O(h^\infty)$.  The artifact measures the finite proxy above after packet analysis through packet-wavefront localization and transport-proxy tests.  Corollary~\ref{cor:packet-wavefront-localization} is therefore the finite wavefront statement used by the branch-identifiability experiments: corrupting the transport branch should enlarge packet-wavefront localization error before it necessarily changes relative $\ell^2$.

\subsection{Restricted H\"ormander-style calculus backbone}
\label{sec:unified-calculus}

The previous subsection treats the hyperbolic case, where the dominant component is a canonical-graph FIO.  The same packet calculus also describes elliptic and parabolic regimes, but the geometry changes.  Elliptic parametrices are pseudodifferential operators with identity canonical relation, while parabolic updates are dissipative pseudodifferential semigroups.  Thus the broader \MiNO\ identity is not ``packets always move.''  It is that the operator is represented by its packet matrix and that the microlocal structure of that matrix is exposed to the architecture.

\begin{theorem}[Restricted FIO/\texorpdfstring{$\Psi$}{Psi}DO packet calculus]
\label{thm:unified-fio-pdo-calculus}
Let $T_h$ be an operator on a compact phase-space window in $\R^d$ admitting the microlocal decomposition
\[
T_h=\sum_{\ell=1}^{L}F_{\ell,h}+A_h+R_h ,
\]
where each $F_{\ell,h}$ is a semiclassical canonical-graph FIO satisfying the off-graph decay and shell-counting hypotheses of Theorem~\ref{thm:packet-sparsity}, $A_h=\Op_h(a)$ is a pseudodifferential operator with $a\in S^m_{1,0}$, and $R_h$ is smoothing of order $R$.  Let $W_h$ be the same packet analysis map used above.  For stencil budgets $K_{\mathrm{fio}}$ and $K_{\mathrm{id}}$, there is a packet operator $\mathcal M_{h,K}^{\mathrm{calc}}$ with
\[
\mathcal M_{h,K}^{\mathrm{calc}}
=
\sum_{\ell=1}^{L}\mathcal M^{\mathrm{fio}}_{\ell,h,K_{\mathrm{fio}}}
+\mathcal M^{\Psi}_{h,K_{\mathrm{id}}}
+\mathcal M^{\mathrm{sm}}_{h},
\]
where the FIO components use canonical tubes, the $\Psi$DO component uses the identity-canonical tube, and the smoothing component is an explicitly bounded residual, such that for every $N>q+1$ and every $M,R\ge 1$,
\begin{align}
\label{eq:unified-calculus-bound}
\|W_hT_hW_h^\ast-\mathcal M_{h,K}^{\mathrm{calc}}\|_{\ell^2\to\ell^2}
\le{}&
C_{\mathrm{fio}}L K_{\mathrm{fio}}^{1-N/q}
+C_{\Psi} K_{\mathrm{id}}^{1-N/q}  \notag\\
&+C_Rh^R
+E_{\mathrm{frame}}(h)+E_{\mathrm{asymp}}(h,M)+C_{\mathrm{cov}}\delta_h .
\end{align}
\end{theorem}

\begin{proof}[Proof sketch]
Apply Theorem~\ref{thm:packet-sparsity} separately to each canonical-graph branch $F_{\ell,h}$.  The sum over finitely many branches contributes the factor $L$.  For $A_h=\Op_h(a)$, use the standard packet almost-diagonalization theorem for pseudodifferential operators: the packet matrix of $A_h$ decays rapidly away from the identity relation $z_\gamma=z_\eta$, with the usual symbol-order envelope inherited from $a\in S^m_{1,0}$.  The same shell-counting and Schur argument as in Theorem~\ref{thm:packet-sparsity} gives the identity-tube truncation term $C_{\Psi}K_{\mathrm{id}}^{1-N/q}$.  Finally, smoothing of order $R$ contributes $C_Rh^R$ on the compact packet window, while finite packetization is recorded by $E_{\mathrm{frame}}$, $E_{\mathrm{asymp}}$, and the normalized covering defect.  Summing the branchwise estimates gives~\eqref{eq:unified-calculus-bound}.
\end{proof}

Theorem~\ref{thm:unified-fio-pdo-calculus} is the restricted H\"ormander-style calculus used by this paper.  It is weaker than full H\"ormander calculus because it assumes a finite branch decomposition, compact phase window, Euclidean packets, and classical packet almost-diagonalization input.  It is stronger than the earlier wave-only bridge because it identifies the three regimes that the model must expose: nontrivial canonical transport, identity-canonical symbol action, and smoothing/dissipative residual structure.

\begin{corollary}[Identity-canonical symbol-order approximation]
\label{cor:identity-symbol-order}
Let $A_h=\Op_h(a)$ with $a\in S^m_{1,0}$ on the compact packet window, and let $\mathcal M^{\Psi}_{h,K,P}$ be a \MiNO\ identity-canonical branch whose order-aware symbol parameterization approximates $a$ in the finitely many symbol seminorms required by Calder\'on--Vaillancourt with error $\epsilon_{s,P}$.  Then, for every Sobolev index $s$ in a fixed compact range,
\begin{align*}
\|W_h^\ast \mathcal M^{\Psi}_{h,K,P} W_h-A_h\|_{H^s\to H^{s-m}}
\le{}&
C_s\bigl(\epsilon_{s,P}+K^{1-N/q}+E_{\mathrm{frame}}(h)+E_{\mathrm{asymp}}(h,M)\\
&\qquad\qquad
+C_{\mathrm{cov}}\delta_h+E_{\mathrm{proj}}\bigr).
\end{align*}
\end{corollary}

\begin{proof}
Classical pseudodifferential calculus gives $A_h:H^s\to H^{s-m}$ with constants controlled by finitely many $S^m_{1,0}$ seminorms of $a$.  Packet almost-diagonalization reduces the operator to an identity-canonical packet matrix, with tail $K^{1-N/q}$ and packetization defects recorded by $E_{\mathrm{frame}}$, $E_{\mathrm{asymp}}$, and the normalized covering term.  Replacing the retained symbol samples by the learned order-aware branch costs $\epsilon_{s,P}$ in the same finite seminorm envelope.  The analysis/synthesis projection contributes $E_{\mathrm{proj}}$.  Combining these estimates gives the displayed Sobolev operator bound.
\end{proof}

Corollary~\ref{cor:identity-symbol-order} states the mathematical target for an order-aware or derivative-regularized symbol branch.  The implementation exposes symbol order and records a finite symbol-order proxy; full H\"ormander-class control would require additional derivative seminorms.

\begin{corollary}[Calculus-backbone \MiNO\ approximation]
\label{cor:calculus-backbone-mino}
Assume the hypotheses of Theorem~\ref{thm:unified-fio-pdo-calculus}.  Let a \MiNO\ block contain matching branch types: bounded canonical FIO branches, an identity-canonical $\Psi$DO symbol branch, a dissipative symbol branch when the target family has parabolic smoothing, and an explicit residual path.  If the learned branch budgets are
\[
\eps_{\mathrm{fio}},\qquad
\eps_{\Psi},\qquad
\eps_{\mathrm{damp}},\qquad
\eps_{\mathrm{res}},
\]
then the packet-space error is bounded by the right-hand side of~\eqref{eq:unified-calculus-bound} plus the corresponding finite learned-budget term
\[
C_{\mathrm{learn}}
\bigl(\eps_{\mathrm{fio}}+\eps_{\Psi}+\eps_{\mathrm{damp}}+\eps_{\mathrm{res}}\bigr).
\]
\end{corollary}

\begin{proof}
Combine the finite branchwise propagation theorem with Theorem~\ref{thm:unified-fio-pdo-calculus}.  The finite part is exactly the machine-checked statement in \codepath{UnifiedCalculus.lean}: FIO-branch error, identity-$\Psi$DO branch error, and smoothing residual error add linearly in packet $\ell^1$.  The continuous theorem supplies the classical truncation and discretization defects.
\end{proof}

\subsection{Admissible microlocal extension axes}
\label{sec:admissible-extension-axes}

The calculus backbone also clarifies which architectural extensions are mathematically admissible.  An extension is admissible when it has (i) a corresponding microlocal object, (ii) an executable model hook, and (iii) a defect term that can be added to the same packet-budget ledger.  This keeps each extension visible to the theorem instead of treating it as an unconstrained residual module.

\begin{table}[t]
\centering
\scriptsize
\setlength{\tabcolsep}{3pt}
\begin{tabular}{p{0.20\linewidth}p{0.25\linewidth}p{0.25\linewidth}p{0.20\linewidth}}
\toprule
Extension axis & Microlocal object & Executable hook & Budget term \\
\midrule
Anisotropic packets / curvelet limit
& Directional packets with parabolic or elongated phase-space tiling
& Current \texttt{anisotropic\_gabor} and \texttt{directional\_gabor} probes; full curvelet/shearlet tokenizer deferred
& $E_{\mathrm{frame}}(\mathfrak C_{\mathrm{frame}},\delta_h,\theta_{\mathrm{dir}})$ \\
Finite multibranch canonical relation
& Finite sum of canonical-graph FIOs on a compact window
& \texttt{num\_canonical\_branches}, metadata or metadata-frequency softmax router, branchwise carrier/landing
& $E_{\mathrm{route}}=E_{\mathrm{prior}}+E_{\mathrm{learned}}$, plus $E_{\mathrm{cross}}+\sum_\ell(E_{\mathrm{tr},\ell}+E_{\mathrm{sym},\ell}+E_{\mathrm{phase},\ell})$ \\
Stronger symbol-class enforcement
& Finite seminorm envelope for $S^m_{1,0}$ and Calder\'on--Vaillancourt bounds
& \texttt{spectral\_order}, identity-$\Psi$DO branch, finite symbol-seminorm proxy
& $E_{S^m}=E_{\mathrm{order}}+E_{\mathrm{seminorm}}$ \\
Outgoing resolvent/radiation symbol
& Local near-characteristic signed resolvent symbol for outgoing high-frequency resolvents
& \texttt{helmholtz\_resolvent} complex-pair multiplier, auto-calibrated characteristic shell, outgoing gate, optional explicit real/imaginary output losses
& $E_{\mathrm{rad}}=E_{\mathrm{shell}}+E_{\mathrm{out}}+E_{\mathbb C}+E_{\mathrm{cap}}+E_{\mathrm{cplx}}$ \\
Nonlinear transport--diffusion coupling
& Linearized quasilinear step: advective FIO composed with dissipative $\Psi$DO
& Transport branch plus dissipative-$\Psi$DO branch and residual pressure/closure path
& $E_{\mathrm{td}}=E_{\mathrm{adv}}+E_{\mathrm{damp}}+E_{\mathrm{split}}+E_{\mathrm{closure}}$ \\
\bottomrule
\end{tabular}
\caption{Microlocally admissible extension axes.  The first four are partially instantiated in the artifact; the last row records the calculus reading of nonlinear or quasilinear transport--diffusion updates and is reserved for a dedicated empirical campaign.}
\label{tab:admissible-extension-axes}
\end{table}

\begin{lemma}[Extension-budget bookkeeping]
\label{lem:extension-budget-bookkeeping}
Assume the hypotheses of Corollary~\ref{cor:calculus-backbone-mino}.  Suppose an architectural extension in Table~\ref{tab:admissible-extension-axes} remains inside the restricted packet calculus and has a verified finite defect $E_{\mathrm{axis}}$ in the appropriate packet operator norm.  Then the same field-level and packet-level conclusions hold with the right-hand side enlarged by $C_{\mathrm{axis}}E_{\mathrm{axis}}$.  For a collection of such extensions,
\[
\|\mathcal M_{\theta,h}^{\mathrm{ext}}-T_h\|
\le
\|\mathcal M_{\theta,h}^{\mathrm{base}}-T_h\|
+C_{\mathrm{frame}}E_{\mathrm{frame}}
+C_{\mathrm{branch}}E_{\mathrm{branch}}
+C_{S^m}E_{S^m}
+C_{\mathrm{rad}}E_{\mathrm{rad}}
+C_{\mathrm{td}}E_{\mathrm{td}} .
\]
Extensions enter the theorem-facing \MiNO-Core claim once their finite defect term is available; otherwise they remain empirical hooks or residual paths.
\end{lemma}

\begin{proof}
Each listed extension changes the retained packet operator by an explicitly named perturbation: frame discretization and anisotropic covering change the analysis/synthesis constants; branch routing and cross-talk perturb the finite FIO sum; symbol enforcement perturbs the finite $S^m$ seminorm envelope; the outgoing resolvent proxy perturbs the near-characteristic shell, the signed radiation multiplier, and its finite complex-pair realization; and transport--diffusion splitting perturbs the product of advective and dissipative packet laws.  The finite bridge is additive in these perturbations, so the result follows by the same analysis--synthesis triangle inequality used in Corollary~\ref{cor:calculus-backbone-mino}.
\end{proof}

\subsection{Machine-checked bridge}

The machine-checked component begins after the classical estimates have organized the packet matrix into canonical shells.  The bridge has a concrete Lean decomposition:
\begin{itemize}
\item \codepath{CanonicalTube.lean} defines the finite shell and canonical-tube objects;
\item \codepath{ShellCount.lean} isolates shell-cardinality and counting interfaces;
\item \codepath{PacketTail.lean} proves the geometric shell-tail estimate;
\item \codepath{AbstractSparsity.lean} packages off-graph decay plus shell counting into an abstract truncation bound;
\item \codepath{RestrictedBridge.lean} packages the restricted non-stationary-phase tail bound with the cross-resolution transfer wrapper;
\item \codepath{StencilPropagation.lean} proves the finite bounded-$K$ propagation theorem used by the continuous-to-discrete stencil corollary;
\item \codepath{StencilWaveTransfer.lean} composes the bounded-stencil coarse theorem with the scale-indexed transfer wrapper;
\item \codepath{WaveTransfer.lean} connects the truncation layer back to the explicit cross-resolution theorem;
\item \codepath{UnifiedCalculus.lean} proves the finite branch-composition spine for FIO, identity-$\Psi$DO, and smoothing residual contributions.
\end{itemize}
In other words, the oscillatory-integral input remains classical, but the shell-counting, truncation, and transfer bridge is now modularized inside the theorem prover.

\subsection{Why the finite core remains \texorpdfstring{$K=1$}{K=1}}
\label{sec:single-packet}

The continuous theory now says that bounded canonical stencils are natural. The Lean-facing finite landing lemma in Section~\ref{sec:theory} nevertheless stays at $K=1$. This is deliberate.

The $K=1$ theorem is the cleanest machine-checkable algebraic core. It isolates transport mismatch, symbol mismatch, and residual error in the smallest exact class. The bounded-$K$ continuous corollary above should therefore be read as an \emph{upgrade path}: short-time wave and abstract FIO sparsity explain why the true operator often lives in a slightly richer class than the exact Lean theorem, while the Lean theorem identifies the primitive that remains visible inside that richer class.

The paper establishes the primitive, exhibits the canonical wave/FIO route to bounded-stencil refinement, and connects the classical continuous analysis to the machine-checked finite bridge.
